\section{The joint center on a single domain}\label{rw:joint-section}

Fix the limiting coefficients
\[
 \bar B=-\bar m=\gamma^2,\qquad
 \bar a=-\bar c=-2\gamma,\qquad \epsilon=\nu=0,
 \qquad P=1/\gamma,\quad\gamma>0.
\]
We use throughout the actual coefficient coordinates of
Proposition~\ref{rw:center-coordinates}:
\[
 \mathbf b=(\bar B,\delta,\tau,\zeta,\nu),\qquad
 \mu_L=\delta-\zeta,\quad\mu_U=\delta+\zeta,\quad
 \tau=\bar a+\bar c.
\]
The lower gap is \eqref{rw:coordinates:lower-gap}. Near $P$, either
gap may be nonpositive, so contact inputs must satisfy both whole-path
passing tests. Beyond $P$, the separated description requires two
positive gaps and continues to $1/p=0$. The difficulty is to control
one comparison through both descriptions as the gaps vanish at unrelated
rates. We use the same original central outputs on their regular overlap
and prove one critical-point direction on their union.

All comparisons below return to the original central height $z$.
Use one fixed outgoing weak reference and one fixed incoming
reference, continued from their fixed inner sections. Write their
complete central maps as $s_++D_+$ and $s_-+J$, and put
\begin{equation}\label{rw:joint:closing}
 F=D_+-J-\sqrt\epsilon\,\nu H_c,\qquad
 K=\epsilon\log(J'/D_+'),\qquad s_--s_+=\sqrt\epsilon\,\nu H_c .
\end{equation}
The prime refers to the same input in the two derivatives. Both
derivatives are nonzero with the same sign on each open passing
interval. A common change of input composes $K$ with that change;
it introduces no extra logarithm. Signed central continuations are
used for coefficient comparison. Positive total central heights and
the actual collar conditions are additionally imposed on counted
orbits.

\begin{proposition}[The joint block]\label{rw:joint-local}
There are a contact neighborhood of $p=P$, a separated interval
continuing to the radial endpoint, a common coefficient neighborhood
and a positive scale cutoff with the following properties.

The two descriptions use the same gaps, phases, central outputs
and coefficient inverse box. Their passing domains here impose the
actual slow passages and fixed regular arcs, while permitting the
signed central comparisons just described. With two positive gaps
these domains have an open overlap and form one interval. With
a nonpositive gap only the retained contact passing interval is used.
No restriction on the ratio of two positive gaps is imposed.

Set $I=|\tau|+|\zeta|$. For one fixed sufficiently small
$\vartheta>0$, every noncenter fibre satisfying
$|\nu|\le\vartheta I$ has at most one zero of $K'$ on this
joint interval, including multiplicity. Each such critical point is
simple and the sign of the second derivative in any regular input
coordinate is $\operatorname{sign}\tau$. Consequently $K$ has
at most two zeros there, with multiplicity.

At an actual closing zero,
$|\nu|\le\eta(\epsilon)I$, where $\eta(\epsilon)\to0$.
Thus this joint block contains at most three isolated counted
cycles, including multiplicity. If $I=0$ and a closing zero exists,
then $\nu=0$ and the comparison is a center identity, contributing
no isolated cycle.
\end{proposition}

Two estimates give the critical-point bound. On the separated interval,
we remove the input-independent action constant and cancel the exact
common powers before forming Hadamard columns. At contact, where division
by a slow factor may be singular, we instead multiply by both factors.
The resulting numerator has bounded jets through signed gaps. Its
limiting value and slope give the same critical-point direction as the
separated columns. A closing zero then supplies the relative coefficient
bound needed to apply these estimates. The estimates include the radial
endpoint; counted finite orbits have strictly positive radial input.

\subsection{Separated expressions before coefficient division}

Choose a small contact radius $r>0$ and then $0<d<r/4$.
The separated finite entries satisfy $p\ge P+d$, and their limiting
radial interval is
\[
 J=[0,(P+d)^{-1}].
\]
Retain slightly larger neighborhoods at every fixed finite boundary.
Place the fixed lower slow sections on opposite sides of the minimum,
with the outer section between $P$ and $P+d$. Choose the upper
sections in the same way. In the upper block first choose their
width for the desired finite normal domination; only afterwards
choose the primary scale for the lower block. These are the
$q=1$ and $q=\epsilon$ choices of
Proposition~\ref{rw:relative-core}.

On a fixed positive annulus in $1/p$, apply
Lemma~\ref{rw:transitions:common-input} before either slow block.
It gives a single coordinate $s_\epsilon$, equal to the actual
radial phase near infinity and to $1/p$ on the finite side. Its
limit on $J$ is $s=1/p$, and its selected transition and inverse
jets are bounded on the overlap. The actual inverse coefficient
chart and these regular input charts are chosen on larger boxes
before the coefficient neighborhood is reduced. Proposition
\ref{rw:domains:separated} then supplies the whole contraction
\[
 (\bar B,\delta,\tau,\zeta,\nu)\longmapsto
 (\bar B,\delta,t\tau,t\zeta,t\nu),\qquad0\le t\le1,
\]
with this common input fixed. Every represented input passes through
both slow blocks along the entire coefficient path. Thus the columns
below compare functions on one input interval, without shrinking that
interval as the coefficients move toward the center.

\begin{lemma}[The nonconstant separated logarithm]
\label{rw:joint:relative-expression}
On this common domain, write $K=C(\mathbf b,\epsilon)+R$ by
subtracting its value at one fixed input. The function $R$ has
uniformly bounded mixed derivatives of total order four, in the
ordinary common input, ordinary coefficients and Euler scale.
The complete deviations have, after possibly decreasing $c>0$,
\begin{equation}\label{rw:joint:absolute}
 \|D_+\|_{C_{s_\epsilon}^1C_{\mathbf b}^1}
       +\|J\|_{C_{s_\epsilon}^1C_{\mathbf b}^1}
       \le Ce^{-c/\epsilon}.
\end{equation}
No coefficient derivative of the subtracted constant is included.
\end{lemma}

\begin{proof}
We treat the finite and radial inputs separately, removing only constants
in the input. On a compact finite patch the first lower regular piece
gives a small normal input
$e^{-S(p)/\epsilon}$, where $S,S_p>c$ with bounded finite jets.
Here the bounded logarithms of regular prefactors have been absorbed
into $S$ with their factor $\epsilon$.
The fixed lower slow map and central connection give
\[
 D_+(p)=C_+L_L e^{-S(p)/\epsilon+r_+(p)}.
\]
Here $C_+L_L$ is input independent. Relative differentiation in
Proposition~\ref{rw:relative-core}, followed by the normalized central
connection, bounds the required jets of $r_+$ by $Ce^{-c/\epsilon}$.
The upper fixed-section map and its regular connectors give an
incoming horizontal phase
\[
 \chi(p)=\alpha\phi(p),\qquad
 \alpha>0,\quad r(p)=\log\phi(p),\quad-r_p>c,
\]
with bounded jets of $r$ and bounded absolute jets of $\alpha$.
The coefficient derivatives of $\log\alpha$ need not be bounded.
The incoming horizontal central map is
$D_I(\chi)=K_I\chi^\mu e^{P_I(\chi)/\epsilon}$, where
\[
 \mu=\frac{1-\epsilon\bar B}{\epsilon\bar B},
 \qquad h_I=\epsilon\mu>c.
\]
The reciprocal preparation in Section~\ref{rw:transitions} supplies
bounded selected jets of the action $P_I$. This form includes the
regular changes of weak coordinate; it does not require bounded
ordinary jets of $P_I/\epsilon$. Normalize $P_I(0)=0$ by changing
$K_I$.
The margin for $-r_p$ follows from its regular normalized upper
entry and the small relative derivative on the chosen normal tube.
Equivalently its zero-scale value is the strictly positive return
rate \eqref{rw:joint:linear-rate} with the opposite sign; the bounded
gap jets extend that margin uniformly to the repeated limit.

Differentiating the actual amplitudes, before taking logarithms,
gives the exact finite-patch expression
\begin{equation}\label{rw:joint:finite-log}
\begin{split}
 K&=C_0+S(p)+h_I r(p)+P_I(\chi)+\epsilon E(p),\\
 E&=-r_+
       +\log(-r_p/S_p)
       +\log(1+\chi P_I'(\chi)/h_I)
       -\log(1-\epsilon r_+'(p)/S_p).
\end{split}
\end{equation}
All logarithm arguments have positive margins. The possibly divergent
actions $\epsilon\log L_L$ and $h_I\log\alpha$ belong to $C_0$.
The nonconstant composition $P_I(\alpha\phi)$ uses absolute jets of
$\alpha\phi$, so its bounds require no derivative of $\log\alpha$.

At the radial end the two horizontal exponents agree exactly for
the same field, by \eqref{rw:transitions:equal-exponents}.
Write the complete outgoing map and the incoming map before their
upper connector as
\[
 D_\pm(t)=K_\pm t^{h_I/\epsilon}
                         e^{P_\pm(t)/\epsilon}.
\]
All slow contractions independent of $t$ belong to $K_\pm$.
The upper connector, including its fixed repeated passage, has
$f(t)=\alpha t e^{\psi(t)}$, where the selected jets of $\psi$
and the absolute jets of $f$ are bounded. Normalize
$P_\pm(0)=\psi(0)=0$ by transferring constants to $K_\pm,\alpha$.
Apply \eqref{rw:transitions:common-power} with $k=\alpha e^\psi$,
grouping $h_I\log\alpha$ into
$C_1=\epsilon\log(K_-/K_+)+h_I\log\alpha$.
The remaining logarithm arguments have positive margins on a fixed
small radius. Its primitive jets follow from the horizontal preparation,
relative slow integrals and central connections. Each block is first
estimated with independent normal input, then composed using absolute
jets of $f$.

Primitive total order five leaves total order four after the one
derivative in \eqref{rw:transitions:common-power} or
\eqref{rw:joint:finite-log}, including $s^3\mathbf b^1$ and
$s^2\mathbf b^2$ after the common input change. The radial and finite
patches form a finite cover of $J$ and describe the same central maps.
Subtracting at a fixed input removes their constants.

Finally each finite lower path contains a fixed positive action
before its slow strip, and the radial input has the common large
power on a fixed interval smaller than one. The fixed inner central
connection contributes a further positive action. The absolute
variational estimates of Section~\ref{rw:transitions} therefore
bound the selected derivatives of both complete deviations by a
finite inverse power of $\epsilon$ times $e^{-c_0/\epsilon}$.
Decrease $c<c_0$ to absorb that power. This proves
\eqref{rw:joint:absolute}, including its continuous radial endpoint.
\end{proof}

\subsection{Actual columns on the whole separated interval}

\begin{proposition}\label{rw:joint:separated}
For $\delta>|\zeta|$, on the common interval $J$,
\begin{equation}\label{rw:joint:Hadamard}
 H:=\partial_{s_\epsilon}K
       =\tau H_\tau+\zeta H_\zeta+\nu H_\nu .
\end{equation}
The three columns are uniformly bounded in $C^2(J)$.
As the scale and coefficient neighborhood tend to the joint limit,
\[
 \|H_\tau-U_0\|_{C^1(J)}
       +\|H_\zeta-V_0\|_{C^1(J)}\longrightarrow0,
\]
where $U_0,V_0$ are \eqref{rw:joint:limit-columns}.
This convergence permits every ratio of the two positive gaps.
\end{proposition}

\begin{proof}
At the center endpoint, Proposition~\ref{rw:domains:separated}
gives $K=H=0$. Thus the Hadamard columns on this same domain are
\[
 H_i(s_\epsilon;\mathbf b)
   =\int_0^1\partial_iH
       (s_\epsilon;\bar B,\delta,t\tau,t\zeta,t\nu)\,\dd t,
 \qquad i=\tau,\zeta,\nu.
\]
Overlapping descriptions give the same integrals. After removing the
constant in $K$, Lemma~\ref{rw:joint:relative-expression} supplies
their $C^2$ bounds from $s^3\mathbf b^1$ and their $C^1$ coefficient
Lipschitz bounds from $s^2\mathbf b^2$, within total order four.

At a fixed positive center gap $g$, the scale limit on every compact
positive subinterval of $J$ is the derivative of
$A_+(p)-A_-(q(p))$. To verify this directly, the lower logarithm
integrates its reduced normal rate to $A_+$ or $A_-$; its relative
logarithms and regular endpoint corrections are multiplied by
$\epsilon$ and tend to zero after input differentiation.
The upper limit is the finite-time return of
\eqref{rw:joint:linear-upper}. At fixed positive gap and finite
$p$, its time and crossings are transverse, so the original upper
equations and their coefficient variational equations converge
there. This proves the limit also after the one coefficient and
one input derivatives used for the columns.

Lemma~\ref{rw:joint:principal-columns} gives that limiting
function in a single expression on the closed interval, with the
exponentially small reciprocal return retained. The finite-scale
columns have a uniform $C^2$ bound. Comparing their values and
first derivatives on $[0,s_2]$ to those at a fixed $s_2>0$,
and then decreasing $s_2$, extends the fixed-gap convergence to
$C^1(J)$, including zero. The limiting formula supplies the same
Lipschitz control there.

The gap derivative bound makes the order of limits explicit. For a center
column
$\mathsf H_\epsilon(g)$ and $0<g\le g_0$, the actual gap
derivative bound gives
\[
 \|\mathsf H_\epsilon(g)-\mathsf H_*\|_{C^1}
 \le Cg_0+
 \|\mathsf H_\epsilon(g_0)-\mathsf H_0(g_0)\|_{C^1}
 +\|\mathsf H_0(g_0)-\mathsf H_*\|_{C^1}.
\]
Choose $g_0$ small first and then the scale. Formula
\eqref{rw:joint:limit-columns} controls the last term, while the
fixed-gap convergence controls the middle term. The same retained mixed
derivatives compare an off-center Hadamard column to its center value
with error
$C(|\tau|+|\zeta|+|\nu|)$ in $C^1(J)$.
All coefficient segments remain in the fixed-mean passing domain.
Thus no estimate depends on
$(\delta+\zeta)/(\delta-\zeta)$.
\end{proof}

\subsection{The contact numerator with two signed gaps}

Use the center-calibrated moving entries of
Proposition~\ref{rw:center-coordinates}, including their fixed
regular initial legs. Denote the entries at the fixed fast sections
by $u(p),v(p)$. They have bounded finite jets and
$u_p,v_p<-c<0$. Their exact slow factors are
\[
 f_L(u)=u^2+\delta-\zeta,\qquad
 f_U(v)=v^2+\delta+\zeta.
\]
Each scalar path ends at a fixed positive exit. A positive gap permits
the whole path; for a nonpositive gap the calibrated entry must exceed
the positive root. Including the finite initial leg gives exactly the
passing test of Proposition~\ref{rw:domains:contact}. Since both entries
decrease with $p$, the two tests select one common interval.

The moving-entry estimates of
Lemma~\ref{rw:transitions:moving-rate}, followed by the incoming
horizontal and central maps, give
\begin{equation}\label{rw:joint:contact-rates}
 \frac{D_{uu}}{D_u}=\frac{A_L(u)}{\epsilon f_L(u)},\qquad
 \frac{J_{vv}}{J_v}=\frac{A_U(v)}{\epsilon f_U(v)}.
\end{equation}
The two numerators have positive margins and bounded finite mixed
jets on the larger auxiliary product domain. More explicitly,
\[
 N_L=\beta_L(u,0)+\epsilon f_L\widetilde j_L',\qquad
 N_U=h_I\beta_U(v,0)+f_U\Psi',
\]
\begin{equation}\label{rw:joint:contact-A}
 A_i=N_i-\epsilon f_i'
                  +\epsilon f_iN_i'/N_i,\qquad i=L,U.
\end{equation}
Here $\widetilde j_L$ is the complete lower relative logarithm.
Writing the upper normal output as
$\chi=C L_Ue^{j_U}$, one may take
$\Psi=h_Ij_U+P_I(\chi)$, with the fixed regular units
included. The absolute derivatives of $\chi$ are bounded even
when its slow exit is lost, so the displayed expression has the
required mixed extension. Choose the contact radius and gaps small
enough that $f_U\Psi'$ does not remove the positive margin of
$h_I\beta_U$. Formula \eqref{rw:joint:contact-A} is the direct
identity obtained from $\partial\log D=N/(\epsilon f)$.
In particular the upper block retains its $q=1$ normal dynamics;
the $\epsilon$ in this logarithmic rate comes from the incoming
horizontal exponent.

For the complete maps $D_+(p)=D(u(p))$, $J_+(p)=J(v(p))$,
differentiate in the unchanged finite input:
\begin{equation}\label{rw:joint:contact-numerator}
\begin{split}
 \mathcal M
 &=f_L(u)f_U(v)K_p\\
 &=f_L(u)A_U(v)v_p-f_U(v)A_L(u)u_p\\
 &\quad+\epsilon f_L(u)f_U(v)
             \left(\frac{v_{pp}}{v_p}-\frac{u_{pp}}{u_p}\right).
\end{split}
\end{equation}
Both multiplying factors are positive on the actual passing interval,
so $\mathcal M$ and $K_p$ have the same zeros there, with the same
multiplicities. The right side extends smoothly to the larger auxiliary
domain even when a passage is absent. Reserve primitive total order
eight here. The two derivatives spent on the logarithmic rates and
connection curvature leave $C_p^3C_{\mathbf b}^3$ bounds for
$\mathcal M$; the ensuing Taylor argument uses only
$C_p^2C_{\mathbf b}^2$. This is a separate budget from the
order-five separated construction.

\begin{proposition}[Contact comparison and transverse jets]
\label{rw:joint:contact}
Let $A=|\tau|+|\zeta|+|\nu|$ and compare with the one center
\[
 \mathbf b_c=(\bar B,\delta+3A,0,0,0).
\]
On the whole original contact passing interval,
\begin{equation}\label{rw:joint:contact-Taylor}
 \mathcal M=\tau C_\tau+\zeta C_\zeta+\nu C_\nu+\mathcal R,
 \qquad \|\mathcal R\|_{C_p^2}\le CA^2,
\end{equation}
where $C_i=\partial_i\mathcal M|_{\mathbf b_c}$ have bounded
$C^2$ norms. At the joint limit the value and first input jets of
$(C_\tau,C_\zeta)$ at $p=P$ are those of
\begin{equation}\label{rw:joint:contact-columns}
 U_*(p)=p^2-p/\gamma,\qquad
 V_*(p)=4p^2/\gamma-2p/\gamma^2
\end{equation}
multiplied by one smooth positive unit. Their Wronskian there is
therefore a positive multiple of $-2/\gamma^4$.
\end{proposition}

\begin{proof}
Proposition~\ref{rw:domains:contact} supplies the entire lifted
coefficient segment and the inclusion of every original passing
input in its slow domain. It also supplies the real center identity
$D_+=J_+$ at its endpoint, for either sign of $\delta+3A$.
Thus $\mathcal M=0$ there. Varying $\delta$ along strictly passing
centers preserves this identity, so its passive $\delta$ derivative
also vanishes, as explained after that proposition. Taylor's formula
on the same larger domain gives \eqref{rw:joint:contact-Taylor}: the
displacement in $\delta$ has length $3A$ but contributes only to the
quadratic remainder. Here $A$ is held fixed for the original fibre;
the absolute values defining it are not differentiated.

We verify the limiting transverse jets before dividing by either
slow factor. In this calculation abbreviate the barred coefficients
by $B,m,a,c$, put $M=-m$, and use
\[
 F_L=1+ap+Bp^2,\qquad F_U=1-cp+Mp^2.
\]
At zero scale and $\nu=0$, the initial lower Morse coordinate is
$-a/2-Bp$, giving the factor $BF_L$. The upper linear field
has the exact scalar coordinate
\[
 v_0=\frac{c/2-MX-cMY/2}{1+MY},\qquad
 v_0'=v_0^2+M-c^2/4,\qquad v_0(p,0)=c/2-Mp.
\]
Its initial factor is $MF_U$. These are the zero-scale limits of
the calibrated entries, up to their finite initial scalar flows.
They use the same strong-fibre normalization fixed in Section
\ref{rw:centers}; they are not independent choices of gap.

The calibrated entries differ from these explicit entries by finite
scalar flows whose times depend on $p$. Retaining this dependence, if
$u_i=\phi^i_{T_i(p)}(e_i(p))$ and
$\lambda_i=\partial_e\phi^i_{T_i}(e_i)>0$, then exactly
\begin{equation}\label{rw:joint:variable-time}
 f_i(u_i)=\lambda_i f_i(e_i),\qquad
 (u_i)_p=\lambda_i(e_{i,p}+f_i(e_i)T_{i,p}).
\end{equation}
The first identity is the scalar variational equation. The times
$T_i$ and their selected derivatives are bounded on these finite
regular legs; the lower scalar time includes its factor $\epsilon$.
After division of \eqref{rw:joint:contact-numerator} by
$\lambda_L\lambda_U$, its first two terms are
\[
 f_L(e_L)A_U(v)e_{U,p}-f_U(e_U)A_L(u)e_{L,p}
 +f_L(e_L)f_U(e_U)(A_U(v)T_{U,p}-A_L(u)T_{L,p}).
\]
The curvature term also retains the product of the two initial
factors. At the joint center each such factor is a positive
multiple of $(1-\gamma p)^2$. A transverse derivative of their
product has zero value and zero first $p$ derivative at $P$.
This explains why variable times and curvature are present in the
full numerator but do not change the two limiting jets.

For a positive upper gap, the exact return rate
\eqref{rw:joint:linear-rate} gives the zero-scale expression
\[
 F_LF_U K_p^0
 =pF_U-pF_L\,
        \frac{M+c/q+1/q^2}{B-a/q+1/q^2}.
\]
The scalar identities \eqref{rw:joint:variable-time} show that
the zero-scale multiplier in \eqref{rw:joint:contact-numerator}
is this product times the positive smooth unit
$BM\lambda_L\lambda_U$. Its long upper transit remains in $q$;
it has not been included among the bounded initial times.

At the joint contact point the reciprocal $1/q$ and the finitely
many mixed derivatives needed for these jets tend to zero.
One way to verify this is the normally dominated moving upper
passage of Section~\ref{rw:transitions}: its normal output is a
positive regular multiple of $1/q$, its transit tends to infinity,
and its differentiated absolute estimates tend to zero. In the
linear field the same conclusion follows from
\[
 e^{cz/\Omega}
 \left(\cos z+\frac{2Mp-c}{\Omega}\sin z\right)=1:
\]
at $p=P$, $z\to\pi/2$ as the gap decreases to zero, and the
reciprocal output decays exponentially in $1/\Omega$.
The mixed extension, rather than a pointwise asymptotic alone,
permits the required transverse differentiation.

Consequently the limiting transverse value and first input jets
are those of the master numerator
\begin{equation}\label{rw:joint:master}
 \mathscr N
 =F_LF_U\left(\frac p{F_L}-\frac{Mp}{BF_U}\right)
 =\frac pB\{\varsigma(1-cp)-M\tau p\},
 \qquad\varsigma=B-M.
\end{equation}
The columns in $(\tau,\varsigma)$ are
$-p^2$ and $p(1-cp)/B$, with Wronskian $p^2/B>0$.
At fixed mean gap the change of transverse coordinates at the
joint center is
$d\varsigma=-\gamma\,d\tau-2d\zeta$.
It gives precisely \eqref{rw:joint:contact-columns}, whose
Wronskian at $P$ is $-2/\gamma^4$.
Multiplication by a smooth positive unit multiplies the determinant
by the square of its value; its derivative only adds a multiple of
the value row to the derivative row.

The positive multiplier can also be identified from the initial
factors. Their difference has transverse part
$Bp\,\tau+(1-cp+2Bp^2)\varsigma$ on the center. It equals
$(-B/p)\mathscr N+2F_L\varsigma$ to first coefficient order.
The last term has zero value and first input derivative at $P$.
The common rate times the decreasing entry derivative is negative,
so its product with $-B/p$ is the same positive unit. This also
checks the sign without dividing by a vanishing factor.

Finally the uniform mixed bounds for the multiplied numerator
transfer these jets to the finite-scale columns and to either sign
of the nearby comparison gap. More explicitly, compare first to a
fixed positive center gap on the same larger domain, use finite
scale convergence there, and then decrease that gap; the retained
coefficient derivative bound controls the first comparison by the
gap difference. This is the same fixed-gap order of choices used
above and does not differentiate a divergent action constant.
\end{proof}

\subsection{One direction across the actual overlap}

\begin{lemma}\label{rw:joint:one-critical}
On the union of the retained contact and separated domains, a
noncenter fibre with $|\nu|\le\vartheta I$, for a common sufficiently
small $\vartheta$, has at most one zero of $K'$, with multiplicity.
At each critical point $|\tau|\ge cI$ and the second derivative
has sign $\tau$ in every regular input coordinate.
\end{lemma}

\begin{proof}
On the separated interval write
$H=\tau U+\zeta V+\nu W$ using
Proposition~\ref{rw:joint:separated}. Its limiting columns satisfy
\[
 V_0<0,\qquad
 (U_0/V_0)'=-\frac{\gamma^2}{2(2\gamma-s)^2}<0
\]
with fixed margins on $J$. The columns have bounded $C^1$ norms
and retain these signs after the restrictions already made.
Lemma~\ref{rw:atlas:two-columns}, with error $\nu W$, gives
$|\tau|\ge cI$ and a simple crossing of sign $\tau$ at every zero.

In contact divide the limiting columns by the positive factor
$p^2$ and their common positive unit. They become
\[
 U=1-\frac1{\gamma p},\qquad
 V=\frac4\gamma-\frac2{\gamma^2p}.
\]
At $P$, $V=2/\gamma>0$ and $(U/V)_p=\gamma^2/2>0$.
The retained higher jets in
Proposition~\ref{rw:joint:contact} preserve these margins on one
fixed smaller contact neighborhood. For the exact Taylor expression
\eqref{rw:joint:contact-Taylor}, its remainder divided by $I$
has $C^1$ norm $O(I)$ when $|\nu|\le\vartheta I$.
The same lemma applies with $C^1$ error $O(|\nu|+I^2)$, giving
$|\tau|\ge cI$ and $\operatorname{sign}\mathcal M_p=\operatorname{sign}\tau$
at each zero. Since $f_Lf_U>0$ on the passing interval,
\eqref{rw:joint:contact-numerator} gives
$\operatorname{sign}K_{pp}=\operatorname{sign}\tau$ there.

These statements concern the same $K$. If $p=p(s_\epsilon)$,
then at a critical point
\[
 K_{s_\epsilon s_\epsilon}=K_{pp}(p_{s_\epsilon})^2.
\]
The positive square, not the sign of the input orientation,
transports its second-derivative sign. Near infinity use the
increasing fibre coordinate $-s_\epsilon$; the same identity
applies to this regular radial coordinate.

There is a genuine open overlap before this direction is used.
Keep the contact radius $r$ already fixed and $0<d<r/4$.
For two positive gaps every entry in a fixed interval, for example
$(P+d,P+2d)$, passes both complete slow paths. It belongs to
the contact neighborhood and to the separated finite patch, with
slightly enlarged intervals at both ends. Their finite initial and
final regular legs are the same actual trajectories by uniqueness,
and their phases and references were fixed together in Section
\ref{rw:centers}. Hence their $D_+,J,K$ agree on this open
interval after the common regular input change.
If one gap is nonpositive, the complete separated strip does not
pass and is not adjoined; the contact whole-path conditions still
select one interval.

The positive-gap union is therefore one interval on which all critical
points are simple and have the same second-derivative sign. Between
two such crossings, $K'$ would have to cross in the opposite direction
or have a nonsimple zero. Both are excluded. Thus the common direction
gives at most one critical point on the whole union.
The argument applies on every compact subinterval, hence on the
whole half-open or compactified interval and at retained finite
endpoint germs.
\end{proof}

\begin{proof}[Proof of Proposition~\ref{rw:joint-local}]
The preparations, domains and critical-point conclusions were proved
above. We supply the forcing and the closing count.
The fixed lower inner exit is on the central side of the repeated
minimum. Its input-independent weak reference therefore continues to
the core independently of that slow passage. The Gaussian calculation
of Lemma~\ref{rw:finite:regular-lower}, on these fixed cores, gives
$s_\pm=\mp\sqrt\epsilon\,\nu H_\pm$ and
$H_c=H_++H_-=\sqrt{2\pi}+o(1)>h_*$ uniformly on the joint box.

On the separated contraction, \eqref{rw:joint:absolute} and its
center identity give
$|D_+-J|\le C(I+|\nu|)e^{-c/\epsilon}$.
The contact absolute maps have the same bound after one coefficient
derivative: their separated entry and fixed exit give a normal
flight of length at least $c/\epsilon$ in the lower block, and
the incoming path has the fixed positive central action.
The absolute estimates and their lost-exit extensions in Section
\ref{rw:transitions} apply. Integrating along the entire lifted
contact path of Proposition~\ref{rw:domains:contact} therefore
gives that same relative bound on every original contact input.
These comparisons concern signed central extensions; they do not
assert positive total heights at all intermediate coefficients.

Lemma~\ref{rw:centers:relative-forcing} consequently gives
$|\nu|\le C\epsilon^{-1/2}e^{-c/\epsilon}I=\eta(\epsilon)I$
on every fixed fibre containing a closing zero.
At a zero the opposite signed references and positive deviations give
positive common height, as in Section~\ref{rw:finite:three-saddle}.

For a noncenter fibre the forcing bound meets the fixed threshold of
Lemma~\ref{rw:joint:one-critical}. Its one-critical-point bound and
Lemma~\ref{rw:zero-calculus} give at most three closing zeros on the
common interval, including multiplicity. Restriction to actual counted
orbits cannot increase this number. For $I=0$ the comparison is the
real center identity. Reflection uses its new outgoing label and inverse
upper passage and preserves Poincare multiplicity by Section~\ref{rw:geometry}.
\end{proof}
